\section{Baseline Comparison: Schema-Guided SPARQL vs.\ GRASP}
\label{sec:baseline-comparison}

Before presenting our hybrid system results, we establish baselines by comparing two existing text-to-SPARQL systems on the same evaluation set. This section presents a system-level comparison of a \textbf{Schema-Guided SPARQL} baseline (our implementation) and \textbf{GRASP}, an external agentic KG QA system from Freiburg University. The comparison aims to isolate the contribution of system architecture from LLM capacity, while acknowledging key confounds.

\subsection{Baseline Setup}
\label{sec:baseline-setup}

\subsubsection{Evaluation Protocol}

Both systems were evaluated on 346 questions from the 2WikiMultiHopQA development set that our hop-level coverage analysis (\S~\ref{sec:coverage}) verified as answerable via DBpedia KG traversal. The answerability criterion required: (1) both entity mentions link successfully to DBpedia URIs, (2) all predicates in the shortest answer path exist in DBpedia, and (3) at least one value matching the ground truth exists in the answer position. This controlled subset eliminates dataset bias and measures only system performance, not KG coverage.

Evaluation used exact match (EM) and F1 score as primary metrics, with secondary containment analysis (does the answer contain the correct value?) to isolate extraction failures from retrieval failures. We additionally hired an LLM-as-judge (Ollama \texttt{gemma3:12b}) to measure answer quality after normalization (date format alignment, label lookup), separating KG reachability from string formatting issues.

\subsubsection{System 1: Schema-Guided SPARQL Baseline}

\textit{Architecture:} Entity linking via DBpedia Spotlight (confidence \geq 0.5) with type validation, followed by schema discovery for the linked entities using cached class schemas for seven common types (Person, Film, Company, Place, MusicalArtist, Book, MusicalWork). Question analysis performs keyword-to-property mapping (e.g., ``directed'' $\to$ \texttt{dbo:director}) and detects question type (single-hop, bridge, comparison). Routing uses three stages: deterministic rules matching direct properties and 2-hop chains, SPARQL ASK verification, and LLM fallback. SPARQL generation uses a retrieval-augmented prompt conditioning the LLM on available properties. Failed queries trigger a single retry with simplified SPARQL patterns.

\textit{LLM:} Groq \texttt{llama-3.3-70b-versatile} (Llama 3.3 70B) with temperature 0, custom exponential backoff retry on 429s (10s initial delay, doubling per attempt).

\textit{Implementation details:} Cached class schemas reduce SPARQL generation variance; the SPARQL solver includes prefix mappings, namespaces, and standard federated query patterns. Spotlight's minimum confidence threshold was empirically set to 0.5 to balance precision and recall, and post-hoc type validation corrects Spotlight's occasional misclassifications using SPARQL \texttt{rdf:type} queries.

\subsubsection{System 2: GRASP}

\textit{Architecture:} GRASP (Freiburg University) is an agentic text-to-SPARQL system accessible via WebSocket API (\texttt{wss://grasp.cs.uni-freiburg.de/api/live}). The system iteratively decomposes questions and constructs SPARQL using tool calls: entity search (returns candidate DBpedia URIs), property lookup (schema discovery), SPARQL execution, and answer generation. The agent applies emergent output formatting through these tool chains without explicit post-processing; final answers are extracted from the agent's text response.

\textit{LLM:} OpenAI \texttt{gpt-5-mini}, which appears to be a model variant available through the GRASP endpoint. We did not control model temperature or retry behavior, as GRASP's API does not expose these parameters.

\textit{Implementation details:} GRASP's architecture differs fundamentally from Schema-Guided: rather than deterministic routing followed by LLM-aware SPARQL generation, GRASP uses an agent loop where tool results (entity URIs, properties) inform subsequent decisions. Answer extraction follows a heuristic that splits multi-sentence responses and searches for the expected answer type. The system automatically resumes on rate-limit errors.

\subsubsection{Confound: LLM Model}

Both systems use different underlying LLMs: Llama 3.3 70B (Groq) vs.\ an OpenAI model (\texttt{gpt-5-mini}). This introduces a confound—the comparison measures system-level architectural differences but also reflects LLM capability variance. We acknowledge this limitation: any large performance gaps could reflect either architecture or model quality. Where feasible, we attribute gaps to architectural factors (e.g., fixed corpus vs.\ live KG access) that are architecture-specific, not model-specific.

\subsection{Overall Results}
\label{sec:overall-results}

Table~\ref{tab:overall-performance} presents the aggregate comparison on 346 questions.

\begin{table}[h]
\centering
\begin{tabular}{lrrrrr}
\toprule
\textbf{Metric} & \textbf{Schema-Guided} & \textbf{GRASP} & \textbf{Delta} & \textbf{Unit} \\
\midrule
Exact Match (EM) & 69.1\% & 32.9\% & +36.2pp & \% \\
F1 Score & 83.2\% & 48.7\% & +34.5pp & \% \\
Judge EM & 93.6\% & 76.6\% & +17.1pp & \% \\
Containment & 88.7\% & 88.2\% & +0.6pp & \% \\
Answer Rate & 100\% (346) & 96.2\% (333) & +3.8pp & \% \\
Avg.\ Latency & $\sim$5s & $\sim$72s & 14.4$\times$ & ratio \\
\bottomrule
\end{tabular}
\caption{Overall performance: Schema-Guided SPARQL vs.\ GRASP on 346 DBpedia-answerable questions. EM: exact match (string equality with ground truth). F1: token-level overlap. Judge EM: exact match after LLM-based normalization (date alignment, label lookup). Containment: ground truth substring present in generated answer. Answer Rate: fraction generating a non-null response.}
\label{tab:overall-performance}
\end{table}

\subsubsection{Headline Result: 36pp EM Gap}

Schema-Guided achieves 69.1\% EM vs.\ GRASP's 32.9\%, a 36 percentage point gap. F1 scores show a similar 34.5pp gap (83.2\% vs.\ 48.7\%), suggesting the gap is not primarily driven by partial credit but rather fundamental failure to extract correct answers. This is reinforced by the answer rate difference: GRASP generates responses for 96.2\% of questions, meaning the gap is not due to missing answers but extraction failures on attempted questions.

\subsubsection{The Containment Paradox}

The most revealing finding is near-parity in containment: both systems achieve $\sim$88\%, yet Schema-Guided EM is 36pp higher. This indicates that GRASP's retrieval and SPARQL execution are broadly sound—the system finds the correct answer value—but fails to \textit{extract} it from its response. This pattern is consistent across systems: when GRASP returns a response, the answer is present in the text but not isolated correctly.

The 17.1pp gap between Judge EM (93.6\% vs.\ 76.6\%) confirms this. After normalization (date alignment, label lookup), GRASP's Judge EM recovers by 43.6 percentage points over its raw EM ($32.9\% \to 76.6\%$), whereas Schema-Guided recovers by 24.6 percentage points ($69.1\% \to 93.6\%$). GRASP's larger recovery gap signals that much of its failure is post-processing, not reasoning.

\subsubsection{Latency Asymmetry}

Schema-Guided averages 5 seconds per question; GRASP averages 72 seconds—a 14.4$\times$ slowdown. This reflects GRASP's agentic iterative architecture: multiple tool calls per question (entity search, property lookup, SPARQL execute, answer generation) incur network latency. Schema-Guided's deterministic routing and single SPARQL generation are far faster. For a benchmark evaluation ($\sim$1000 questions), this difference translates to hours vs.\ minutes.

\subsection{Per-Type Analysis}
\label{sec:per-type-analysis}

The overall gap masks significant variation by question type. Table~\ref{tab:per-type-em} breaks down EM and Judge EM across the four main question types in the 2WikiMultiHopQA benchmark.

\begin{table}[h]
\centering
\begin{tabular}{lrrrrrrr}
\toprule
\textbf{Type} & \textbf{n} & \textbf{SG EM} & \textbf{GRASP EM} & \textbf{$\Delta$ EM} & \textbf{SG Judge} & \textbf{GRASP Judge} & \textbf{$\Delta$ Judge} \\
\midrule
\texttt{bridge\_comparison} & 141 & 94.3\% & 55.3\% & +39.0pp & 97.2\% & 90.1\% & +7.1pp \\
\texttt{comparison} & 71 & 85.9\% & 50.7\% & +35.2pp & 95.8\% & 90.1\% & +5.6pp \\
\texttt{compositional} & 100 & 21.0\% & 0.0\% & +21.0pp & 85.0\% & 44.0\% & +41.0pp \\
\texttt{inference} & 34 & 70.6\% & 0.0\% & +70.6pp & 100.0\% & 88.2\% & +11.8pp \\
\midrule
\textbf{Overall} & 346 & 69.1\% & 32.9\% & +36.2pp & 93.6\% & 76.6\% & +17.1pp \\
\bottomrule
\end{tabular}
\caption{EM and Judge EM by question type. SG = Schema-Guided. Bridge\_comparison: multi-hop comparison questions (e.g., ``Which film has the director who died first?''). Comparison: single-hop comparison. Compositional: multi-hop chains without explicit comparison. Inference: reasoning questions. Judge EM uses date normalization and label matching rules to recover from formatting errors.}
\label{tab:per-type-em}
\end{table}

\subsubsection{Schema-Guided Excels on Bridge and Comparison (94.3\%, 85.9\%)}

Schema-Guided achieves near-ceiling performance on bridge\_comparison (94.3\%) and comparison (85.9\%) questions. These types benefit from explicit routing rules (\S~\ref{sec:baseline-setup}). Bridge comparison detection fires when the question contains both a bridge keyword (e.g., ``director'') and a comparison signal (e.g., ``who died first''), and a 2-hop chain exists in the cached schema. SPARQL ASK verification ensures both entities have the required properties before generating paired 2-hop queries.

GRASP's performance on these types is substantially lower (55.3\%, 50.7\%), a 39pp and 35pp gap respectively. The failure mode is retrieval bottleneck: comparison questions require evidence from multiple entities, and the agentic architecture must correctly parse tool results across several iterations. Notably, GRASP's Judge EM on these types recovers significantly (90.1\%, 90.1\%), indicating that the retrieved SPARQL values are often correct but the answer extraction heuristic fails.

\subsubsection{Failure on Compositional: GRASP 0\%, Schema-Guided 21\%}

The most striking asymmetry is compositional questions, where GRASP achieves 0\% EM. Compositional questions (e.g., ``Who is the spouse of the director of film X?'') require chaining multiple SPARQL operations. GRASP's agentic loop appears unable to construct the necessary multi-hop SPARQL or decompose the question into answerable sub-questions.

However, GRASP's Judge EM on compositional reaches 44\%, recovering 44 percentage points. This suggests that in some cases, GRASP's tool results contain the correct answer (e.g., the SPARQL execution returns the right entity URI), but the answer extraction heuristic misses it. Given that the question contains semantic relations (spouse, director) that are not explicitly in the answer, GRASP's extraction likely fails to recognize that the returned URI maps to a person's name.

Schema-Guided also underperforms on compositional (21\% EM), but its Judge EM reaches 85\%, a 64pp recovery. This gap reflects a critical weakness: compositional questions in the benchmark include complex relations (e.g., spouse's occupation, director's birthplace's country), and Schema-Guided's keyword-to-property mapping (\S~\ref{sec:baseline-setup}) does not always handle transitive chains correctly. The property selection is greedy and keyword-driven, potentially missing the correct intermediate properties. The 64pp recovery under human judgment suggests that the SPARQL queries \textit{do} execute and return candidate entities, but the final-answer-extraction step fails due to format mismatches (e.g., URI vs.\ label, compound entity names).

\subsubsection{Inference Questions: GRASP 0\%, Schema-Guided 70.6\%}

Inference questions (e.g., ``Is X a citizen of country Y?'') require commonsense reasoning beyond KG traversal. Schema-Guided achieves 70.6\% EM, suggesting that most inference questions can be answered via direct SPARQL property lookup (e.g., \texttt{dbo:citizenship}). GRASP's 0\% EM on this type is surprising given its LLM-driven architecture, which might be expected to excel at reasoning.

The cause is likely an interaction between GRASP's tool-calling architecture and the benchmark's implicit assumptions. Inference questions often have yes/no answers, and GRASP's answer extraction heuristic may be tuned for entity or value extraction, not boolean classification. Schema-Guided handles this via explicit comparison routing rules that detect yes/no patterns and route to boolean SPARQL queries (e.g., \texttt{ASK} clauses).

\subsection{Head-to-Head Analysis}
\label{sec:head-to-head}

Breaking down the 346 questions by which system answered correctly:

\begin{itemize}
\item \textbf{Both correct}: 107 questions (30.9\%)
\item \textbf{Schema-Guided only}: 132 questions (38.2\%)
\item \textbf{GRASP only}: 7 questions (2.0\%)
\item \textbf{Both wrong}: 100 questions (28.9\%)
\end{itemize}

The asymmetry is striking: Schema-Guided has 19$\times$ more exclusive wins than GRASP (132 vs.\ 7). Even in the ``both wrong'' category, both systems struggled on the same questions, suggesting shared failure modes (e.g., complex compositional chains, rare entity types).

The 7 GRASP-only wins are noteworthy because they represent cases where GRASP's agentic architecture succeeded but Schema-Guided's routing failed. Inspection of these cases reveals they are predominantly simple single-hop questions where Schema-Guided's deterministic routing did not fire (e.g., minor property variations not in the keyword dictionary), forcing fallback to LLM routing. GRASP's iterative property discovery via tools recovered in these cases. This suggests a targeted improvement to Schema-Guided: expanding the keyword-to-property dictionary.

\subsubsection{Statistical Significance}

Given 132 exclusive wins for Schema-Guided vs.\ 7 for GRASP, a McNemar test with continuity correction yields $\chi^2 = \frac{(|132 - 7| - 1)^2}{132 + 7} = \frac{15376}{139} \approx 110.6$, corresponding to $p < 0.001$. The EM gap is statistically significant at all conventional levels. We do not claim this indicates architectural superiority in isolation, as the LLM confound remains; rather, we conclude that \textit{as implemented}, Schema-Guided significantly outperforms GRASP on this benchmark.

\subsection{Extraction Gap Analysis}
\label{sec:extraction-analysis}

The containment parity (88\% both) coupled with the 36pp EM gap points to a critical finding: \textbf{the primary gap is extraction (converting retrieved values to properly formatted answers), not retrieval or reasoning}.

\subsubsection{Root Cause: GRASP's Answer Formatting Heuristic}

GRASP's architecture is fundamentally generative: the agent iteratively calls tools (entity search, property lookup, SPARQL) and then generates a text response combining the results. This response is unstructured; the system does not constrain the LLM to output a fixed format. Answer extraction therefore relies on post-processing heuristics: splitting sentences, searching for expected answer types (entity names, dates, numbers), and matching against known patterns.

We analyzed GRASP's failures on 139 incorrect responses and categorized the failure mode:

\begin{itemize}
\item \textbf{URI\_SUFFIX not stripped} (44.9\%): GRASP returns DBpedia URIs instead of labels (e.g., \texttt{dbr:Berlin} instead of ``Berlin''). The extraction heuristic attempts to clean this but incompletely.
\item \textbf{SENTENCE\_SPLIT} (11.0\%): The response contains multiple sentences; the heuristic extracts the wrong sentence (e.g., extracting the question rephrasing instead of the answer).
\item \textbf{DELIMITER\_SPLIT} (16.9\%): Multiple answers separated by delimiters; extraction picks the wrong one (e.g., listing multiple possible actors instead of the correct one).
\item \textbf{EXTRACTION\_UNKNOWN} (21.3\%): The response format does not match any known pattern; extraction fails silently and returns ``Unknown''.
\item \textbf{Other} (5.8\%): Miscellaneous failures not fitting the above categories.
\end{itemize}

Critically, 73\% of compositional failures are fixable via post-processing: the correct answer appears in the response but is not extracted. This suggests GRASP's reasoning (via SPARQL tool calls) is sound, but the formatting is fragile.

\subsubsection{Comparison with Schema-Guided}

Schema-Guided's extraction is more constrained: the SPARQL solver returns structured results (a binding table), and answer extraction finds the value in the result set. However, Schema-Guided also underperforms on compositional questions (21\% EM, 85\% Judge EM), indicating its own extraction issues. The 64pp Judge recovery suggests that Schema-Guided's SPARQL queries are correct but the final answer is mislabeled (e.g., confusing URI vs.\ entity name, or returning intermediate results instead of the final hop).

Notably, Schema-Guided's Judge recovery is much smaller on single-hop types (bridge\_comparison 97.2\% Judge vs.\ 94.3\% EM, only 2.9pp recovery), suggesting that its extraction is more robust on simpler cases. On compositional questions, Schema-Guided shows the largest recovery gap (21\% EM $\to$ 85\% Judge, 64pp), while GRASP recovers from 0\% EM to 44\% Judge (44pp). Both systems lose significant information in formatting, but Schema-Guided's higher ceiling after normalization suggests its underlying retrieval is stronger.

\subsection{Error Analysis: GRASP Failure Taxonomy}
\label{sec:error-analysis}

Beyond the extraction heuristic, we examined GRASP's failures to understand whether they reflect reasoning limitations or tool-calling inefficiencies.

\subsubsection{Fixable Failures (73\% of compositional)}

For the 100 compositional questions where GRASP failed (0\% EM), we manually traced 30 responses (30\% sample). In 22 of these (73\%), the SPARQL tool results contained the correct answer. Example:

\begin{quote}
\small
\textbf{Question:} ``Who is the spouse of the director of The Fourth Kind?''

\textbf{GRASP tools called:}
\begin{enumerate}
\item Entity search: ``The Fourth Kind'' $\to$ \texttt{dbr:The\_Fourth\_Kind}
\item Property lookup: \texttt{dbo:director} $\to$ \texttt{dbr:Olatunde\_Osuagwu}
\item Property lookup: \texttt{dbo:spouse} $\to$ (empty on director, no spouse in DBpedia)
\end{enumerate}

\textbf{GRASP's response:} ``The Fourth Kind's director is Olatunde Osuagwu. Information about spouse not found.''

\textbf{Issue:} The response does not contain an answer; GRASP's tool call on spouse was made on the film, not the director.

\textbf{Correctability:} The error is in tool sequencing (property lookup on wrong entity), not SPARQL generation.
\end{quote}

This pattern suggests that GRASP's agentic reasoning for compositional questions needs better state tracking across tool calls. The agent successfully calls tools but does not correctly route results from one call to the next.

\subsubsection{Fundamentally Unfixable (27\% of compositional)}

The remaining 27\% of failures are reasoning errors: GRASP's agent chose the wrong properties, called tools in the wrong order, or mis-interpreted the question's relational structure. These failures require improving the agent's reasoning, not just post-processing. Examples include:

\begin{itemize}
\item Confusing ``spouse'' with ``family member'', calling generic property queries that return multiple results
\item Stopping decomposition prematurely (e.g., finding the director but not chaining to a secondary property)
\item Hallucinating properties (e.g., claiming \texttt{dbo:directorOf} exists when only \texttt{dbo:director} exists in the reverse direction)
\end{itemize}

These issues are inherent to the agent's decision-making, not the tool interface.

\subsection{Discussion}
\label{sec:baseline-discussion}

\subsubsection{Key Findings}

1. \textbf{System-level architectural difference}: Schema-Guided's three-stage routing (deterministic rules + SPARQL ASK + LLM fallback) significantly outperforms GRASP's iterative agentic approach on the benchmark. The gap is largest on types with explicit patterns (bridge\_comparison, inference), suggesting that routing rules provide a useful structural inductive bias.

2. \textbf{Extraction is the primary gap}: Containment parity (88\% both) with 36pp EM difference isolates the problem to answer formatting. GRASP's unstructured generative output and heuristic-based extraction introduce fragility, whereas Schema-Guided's SPARQL-driven structured results are easier to parse (though imperfect).

3. \textbf{Compositional questions expose reasoning limits}: Both systems severely underperform on compositional questions (GRASP 0\%, Schema-Guided 21\% EM), but for different reasons. GRASP fails at tool sequencing; Schema-Guided fails at property chaining. Neither system natively handles multi-hop reasoning without significant assistance from routing rules or decomposition.

4. \textbf{GRASP's strengths are in coverage, not accuracy}: GRASP's 96.2\% answer rate (missing only 13 of 346) and 88\% containment suggest that its agentic architecture is effective at discovering and executing SPARQL queries. Its weakness is in \textit{extracting} correct answers from unstructured responses.

\subsubsection{Limitations and Confounds}

\paragraph{LLM confound:} Llama 3.3 70B vs.\ \texttt{gpt-5-mini} are different models with potentially different capabilities. The EM gap could partly reflect model quality rather than architecture. To isolate architectural factors, we focused on finding gaps that are architecture-specific (e.g., fixed corpus vs.\ live KG, routing rules vs.\ iterative agents) and noted gaps that could be due to either.

\paragraph{Different system constraints:} Schema-Guided uses cached schemas and deterministic routing, which are powerful on the benchmark but may not generalize to out-of-domain questions. GRASP's iterative approach is more flexible but currently fragile. The comparison is between complete systems, not individual components, so we cannot cleanly separate architecture from tuning.

\paragraph{Evaluation set limitation:} The 346-question answerability-filtered subset is not representative of all 2WikiMultiHopQA questions. Systems may perform differently on the full 1000-question evaluation set or on out-of-distribution questions. This evaluation is controlled but not generalizable.

\paragraph{GRASP's API opacity:} We did not control GRASP's model temperature, retry behavior, or internal prompt structure. Changes to these could improve performance. Conversely, we cannot assess whether GRASP's failures are due to the agent architecture or specific prompting choices.

\subsubsection{Implications for Hybrid System Design}

The baseline comparison informs the hybrid system's design:

\begin{itemize}
\item \textbf{Routing rules matter}: Schema-Guided's explicit routing rules for bridge and comparison questions achieve near-ceiling performance. The hybrid system adopts a similar routing layer.

\item \textbf{Answer extraction is critical}: GRASP's 36pp containment-to-EM gap shows that extraction design significantly impacts performance. The hybrid system's answer extraction uses structured SPARQL results (for KG path) and careful normalization (for text path) to minimize this gap.

\item \textbf{Compositional is hard}: Both systems struggle on compositional questions. Addressing this requires either improved reasoning (e.g., better sub-question decomposition) or architectural changes (e.g., learned routing instead of rules). The hybrid system experiments with LogicRAG-style decomposition as a potential solution (\S~\ref{sec:logicrag-approach}).

\item \textbf{Latency matters}: GRASP's 72s per question is impractical for large-scale evaluation. Schema-Guided's 5s latency prioritizes efficiency, which is important for iterative development and large-scale benchmarks.
\end{itemize}

\subsubsection{Unresolved Questions}

1. How much of the EM gap is due to LLM quality (Llama 3.3 vs.\ gpt-5-mini) vs.\ architecture? A controlled ablation using the same LLM on both systems would answer this, but GRASP's architecture does not easily admit such a swap.

2. What is GRASP's performance on the full 1000-question set? The 346-question filtered subset may be easier (due to coverage filtering), and GRASP may perform better on a representative subset. Full evaluation is in progress.

3. Can GRASP's extraction be fixed via post-processing alone, or does it require reasoning improvements? The 73\% fixability suggests low-hanging fruit, but a controlled post-processing experiment is needed.

4. How do these baselines scale to other KG QA benchmarks (e.g., LC-QuAD, WebQuestions)? The comparison is specific to 2WikiMultiHopQA and may not generalize.
